%==============================================================
% BAB V TRANSPARANSI PENGGUNAAN KECERDASAN BUATAN
%==============================================================

\section{Transparansi Penggunaan Kecerdasan Buatan}

\subsection{Pernyataan Transparansi AI dalam Siklus Pengembangan}

Sejalan dengan standar integritas akademik dan arahan Kementerian Kesehatan RI mengenai etika kecerdasan buatan, tim pengembang menyampaikan pernyataan terbuka terkait pemanfaatan alat bantu AI generatif dalam proses perancangan, penulisan proposal, dan rekayasa TropiCare. AI diposisikan secara ketat sebagai akselerator kognitif, penyunting tata bahasa, dan asisten kode (\textit{coding copilot}), tanpa menggantikan orisinalitas pemikiran, pertimbangan klinis, maupun tanggung jawab tim manusia (Tabel~\ref{tab:transparansi_ai}).

\begin{table}[H]
\centering
\small
\caption{Matriks Transparansi Pemanfaatan AI dalam Siklus Proyek}\label{tab:transparansi_ai}
\vspace{3pt}
\begin{tabularx}{\textwidth}{p{3.2cm} X X}
\toprule
\textbf{Fase Aktivitas} & \textbf{Kontribusi Alat Bantu AI} & \textbf{Verifikasi \& Pengawasan Manusia} \\
\midrule
\textbf{Perancangan Ide \& Konsep} & Eksplorasi skenario risiko (\textit{edge-cases}), simulasi variasi kosakata awam untuk uji ketahanan leksikal, dan komparasi literatur sistem triase global. & Seleksi kritis oleh tim berdasarkan realitas faskes primer Indonesia dan keselarasan dengan regulasi Kemenkes RI. \\
\midrule
\textbf{Penyusunan Proposal} & Perumusan draf awal kalimat, perbaikan tata bahasa formal, dan penataan struktur matriks komparasi. & Penyuntingan substansial, verifikasi silang dokumen regulasi (UU PDP, UU Kesehatan, KMK PPK FKTP), dan penulisan ulang analisis kausal oleh tim. \\
\midrule
\textbf{Rekayasa Kode \& Desain} & Optimasi sintaksis dokumen LaTeX/TikZ serta generasi kerangka kode dasar (\textit{boilerplate}) FastAPI dan React. & Penulisan mandiri logika inti (algoritma pencocokan KMK, \textit{rule-engine} 11 \textit{Red Flags}, serialisasi QR), tinjauan kode baris-demi-baris, dan pengujian unit. \\
\bottomrule
\end{tabularx}
\end{table}

\subsection{Batasan Sistem dan Pengawasan Manusia (\textit{Human-in-the-Loop})}

Tim TropiCare menegaskan bahwa kecerdasan buatan hanyalah alat bantu (\textit{assistive tool}), bukan pengambil keputusan klinis otonom. Model AI bekerja berbasis probabilitas statistik dan rentan halusinasi jika tidak dibatasi. Oleh karena itu, TropiCare menerapkan tiga pilar pengawasan manusia:

\begin{enumerate}
    \item \textbf{Validasi Medis oleh Kolaborator Klinis:} Seluruh logika pohon keputusan klinis, pembagian 20 penyakit tropis, penetapan 11 kriteria tanda bahaya (\textit{Red Flags}), dan alur rujukan dikurasi langsung oleh kolaborator mahasiswa kedokteran berdasarkan KMK No. HK.01.07/MENKES/1186/2022 \citep{Kemenkes2022a} dan KMK No. HK.01.07/MENKES/1936/2022 \citep{Kemenkes2022b}. Hal ini menjamin kepatuhan pada Standar Kompetensi Dokter Indonesia \citep{KKI2012}.
    \item \textbf{Rekayasa Perangkat Lunak Deterministik:} Arsitektur backend FastAPI, frontend React, basis data direktori 13.930 faskes SATUSEHAT MSI, dan pencarian faskes terdekat dirancang mandiri oleh mahasiswa ilmu komputer tanpa ketergantungan logika otonom pihak ketiga.
    \item \textbf{Keputusan Etis Pembatasan Fitur:} Tim secara sadar menghapus fitur kamera ruam kulit demi mematuhi UU Pelindungan Data Pribadi No. 27/2022 \citep{UUPDP2022} mengenai data biometrik sensitif dan mencegah bias warna kulit sawo matang/gelap, serta menolak penyimpanan nomor KTP/NIK di server demi keamanan privasi masyarakat.
\end{enumerate}

\subsection{Keselarasan dengan Panduan Etika AI Kemenkes RI dan UNESCO}

Operasionalisasi TropiCare berpedoman teguh pada empat pilar etika kecerdasan buatan Kemenkes RI \citep{KemenkesDTO2024} dan rekomendasi UNESCO \citep{UNESCO2021} (Tabel~\ref{tab:etika_ai}).

\begin{table}[H]
\centering
\small
\caption{Keselarasan Implementasi TropiCare dengan Prinsip Etika AI Nasional dan Global}\label{tab:etika_ai}
\vspace{3pt}
\begin{tabularx}{\textwidth}{p{3.5cm} X X}
\toprule
\textbf{Prinsip Etika} & \textbf{Mandat Regulasi} & \textbf{Implementasi Konkret pada TropiCare} \\
\midrule
\textbf{Transparansi \& Keterjelasan (\textit{Explainability})} & AI medis tidak boleh bersifat kotak hitam (\textit{black-box}); penalaran sistem harus dapat diaudit. & Fitur \textbf{``Kenapa aku menebak ini?''} yang menampilkan kutipan langsung kriteria klinis KMK PPK FKTP 2022 sebagai dasar rekomendasi. \\
\midrule
\textbf{Akuntabilitas Mediko-Legal} & Batasan tegas fungsi perangkat lunak dengan wewenang diagnosis dokter berlisensi. & Sistem dibatasi murni pada \textbf{triase pra-hospital non-diagnosis}; sistem tidak pernah meresepkan antibiotik mandiri (mencegah resistensi kuman). \\
\midrule
\textbf{Keamanan Pasien (\textit{Safety-First})} & Pencegahan keterlambatan penanganan akibat kegagalan komputasi. & Deteksi 11 \textit{Red Flags} di peramban yang melompat langsung ke rute darurat IGD/119 dalam $<500\text{ ms}$ tanpa menunggu server. \\
\midrule
\textbf{Pelindungan Privasi (\textit{Privacy-by-Design})} & Penghormatan kedaulatan individu atas data kesehatan pribadi. & Arsitektur \textbf{Zero-NIK}, kalkulasi GPS lokal tanpa disimpan di server, dan pertukaran data klinis via QR Code luring. \\
\bottomrule
\end{tabularx}
\end{table}
